\documentclass{ibetest}
\begin{document}
\maketest{Progress Test 3.A --- Thermo-S1}
         {3.A \,$\cdot$\, Low-density gases: the historical laws}
         {7 minutes}{14}

% ---------------------------------------------------------------------------
\exercise{Unit conversion}{1 mark}
The product $pV$ of a gas equals $\qty{2.0}{bar.L}$. Express it in joules.
\answer{1}{2.2cm}{%
  $pV = 2.0\times10^{5}\ \mathrm{Pa}\times1.0\times10^{-3}\ \mathrm{m^3}$ \qquad
  $\boxed{pV = 2.0\times10^{2}\ \mathrm J}$ \quad ($\mathrm{Pa\,m^3} = \mathrm{N\,m} = \mathrm J$)}
\begin{scheme}
\pt{1}{1 pt}{$200$~J, with the unit.}
\end{scheme}
\begin{tip}
$pV$ has the dimension of an energy, which is why $R$ is in $\mathrm{J\,K^{-1}\,mol^{-1}}$.
Useful to remember: $1\ \mathrm{bar\,L} = 100\ \mathrm J$.
\end{tip}

% ---------------------------------------------------------------------------
\exercise{Boyle--Mariotte law}{5 marks}
State the Boyle--Mariotte law: the conditions under which it holds, and its formula. Deduce
from it the isothermal compressibility $\chi_T$ of a low-density gas.
\answer{2,2,1}{6.2cm}{%
  \textbf{Boyle--Mariotte (1662):} for a low-density gas at fixed temperature $T$ (and fixed
  $n$), the product $pV$ is constant, whatever the gas:\qquad
  $\boxed{pV = \text{const.}\quad\text{or}\quad p_1V_1 = p_2V_2\ \ (T_1 = T_2)}$\\[5pt]
  So $V = C/p$ and
  $\left(\dfrac{\partial V}{\partial p}\right)_{\!T} = -\dfrac{C}{p^2} = -\dfrac Vp$, hence\qquad
  $\boxed{\chi_T = -\dfrac1V\left(\dfrac{\partial V}{\partial p}\right)_{\!T} = \dfrac1p}$
  \ in $\mathrm{Pa^{-1}}$.\\[3pt]
  The gas becomes less compressible as the pressure increases.}
\begin{scheme}
\pt{1}{2 pts}{Statement: fixed $T$ and $n$, low density (1~pt); $pV = $ const.\ (1~pt).}
\pt{2}{2 pts}{$\left(\partial V/\partial p\right)_T = -V/p$, derived.}
\pt{3}{1 pt}{$\chi_T = 1/p$, in $\mathrm{Pa^{-1}}$.}
\end{scheme}

% ---------------------------------------------------------------------------
\exercise{The three historical laws}{3 marks}
\question{Charles's law holds at fixed:}
\opts{0}{$n$ and $T$}{1}{$n$ and $p$}{0}{$n$ and $V$}{0}{$p$ and $V$}
\why{Tab.~4: Charles (1787), fixed $n$ and $p$, free $V$ and $T$: $V/T = $ const.}
\question{Gay-Lussac's law states that, at fixed $n$ and $V$:}
\opts{0}{$pV = $ const.}{0}{$V/T = $ const.}{1}{$p/T = $ const.}{0}{$pT = $ const.}
\why{equation~(14) of the notes: $p_1/T_1 = p_2/T_2$ if $V_1 = V_2$.}
\question{For a low-density gas, the isobaric expansion coefficient is:}
\opts{1}{$\alpha = 1/T$}{0}{$\alpha = 1/p$}{0}{$\alpha = T$}{0}{$\alpha = 0$}
\why{$V/T = $ const.\ gives $(\partial V/\partial T)_p = V/T$, so $\alpha = 1/T$ (Box~15).}

% ---------------------------------------------------------------------------
\exercise{A rising bubble and a heated balloon}{5 marks}
(a) An air bubble of volume $V_1$ is released at a depth where the pressure is $p_1$. It rises
to the surface, where the pressure is $p_2$, at constant temperature. Find its volume $V_2$ at
the surface.
(b) A balloon of volume $V$ at $\theta_1$ is heated to $\theta_2$ at constant pressure. Find
its new volume $V'$.
\data{(a) \ $V_1 = \qty{2.0}{\centi\metre\cubed}$ \hspace{8mm} $p_1 = \qty{3.0}{bar}$ \hspace{8mm}
      $p_2 = \qty{1.0}{bar}$\\[2pt]
      (b) \ $V = \qty{3.0}{L}$ \hspace{8mm} $\theta_1 = \qty{27}{\degreeCelsius}$ \hspace{8mm}
      $\theta_2 = \qty{77}{\degreeCelsius}$}
\answer{1,1,1,1,1}{6.2cm}{%
  (a) Fixed $n$ and $T$, so Boyle--Mariotte applies: $p_1V_1 = p_2V_2$, i.e.
  $\boxed{V_2 = \dfrac{p_1V_1}{p_2}} = \dfrac{3.0}{1.0}\times2.0\ \mathrm{cm^3}$ \qquad
  $\boxed{V_2 = 6.0\ \mathrm{cm^3}}$\\[5pt]
  (b) Fixed $n$ and $p$, so Charles applies: $\dfrac{V}{T_1} = \dfrac{V'}{T_2}$, i.e.
  $\boxed{V' = V\,\dfrac{T_2}{T_1}}$, with \textbf{absolute} temperatures
  $T_1 = 300\ \mathrm K$ and $T_2 = 350\ \mathrm K$:\\[3pt]
  \hspace*{5mm}$V' = 3.0\ \mathrm L\times\dfrac{350}{300}$ \qquad $\boxed{V' = 3.5\ \mathrm L}$}
\begin{scheme}
\pt{1}{1 pt}{$V_2 = p_1V_1/p_2$ (literal).}
\pt{2}{1 pt}{$V_2 = 6.0\ \mathrm{cm^3}$.}
\pt{3}{1 pt}{$V' = VT_2/T_1$ (literal).}
\pt{4}{1 pt}{Temperatures converted into kelvins: $300$~K and $350$~K.}
\pt{5}{1 pt}{$V' = 3.5$~L.}
\end{scheme}
\begin{tip}
The laws of Charles and Gay-Lussac need absolute temperatures. Using $^\circ$C would give
$V' = 3.0\times77/27\approx8.6$~L, which is absurd.
\end{tip}

\finish
\end{document}
